\documentclass[a4paper,11pt]{article}

\usepackage{graphicx}
\usepackage{amsmath,amssymb}
\usepackage{hyperref}
\usepackage[a4paper,margin=2.5cm]{geometry}

\title{INF-203 Project Report}
\author{''}
\date{June 2026}

\begin{document}

\maketitle

\tableofcontents
\newpage

\section{Introduction}
Brief introduction to the project and objectives.

\section{Task 1}
Our group created a shared GitHub repository for the project. The repository is used to store code and project files, allowing all members to contribute and track changes through version control. This setup helps document each member's contributions and supports collaborative development throughout the course.
The repository can be accessed here: \href{https://github.com/izzaqamar/Inf-203.git}{Inf-203 GitHub repository}
\section{Task 2}

We created a requirements register using a spreadsheet to manage the system requirements. The requirements are structured using an Agile approach with epics and user stories, and are prioritised using the MoSCoW method.

Two main epics were defined: \textit{Data Processing and Output} and \textit{User Feedback and Support}. Each epic contains user stories describing system functionality from different user perspectives. The requirements include three personas: Student, Teaching Assistant, and Teacher.
The requirements register will be maintained throughout the project and updated as the system is developed and refined.
The report can be accessed here: \href{https://docs.google.com/spreadsheets/d/187aslKKx-FqcoSc7nYzcPrt7RW5b8iT2UYNy3XJ90Yw/edit?usp=sharing}{Inf-203 Latex Report}
\section{Task 3}

\section{Task 4}
We created a function that loads a jsonld file, extracts the triples, and stores them in a list. The function then uses the list to create a graph using the triples, before it returns the graph.
Additionally we created a quick program that visualizes the graph using the Kamada-Kawai layout. The graph is visualized with nodes representing subjects and objects, and edges representing predicates. 

\section{Task 5}

\section{Task 6}

\section{Task 7}

\section{Task 8}

\section{Task 9}
We implemented this by extending the query_path method to accept a directions list alongside the labels list, and an optional filters list:

def query_path(self, labels, directions, filters=None):

Each entry in directions is a boolean where True means follow the edge forward and False means follow it backward. The filters list accepts callable filter objects, where all filters must pass for a node to be included in the path. This allows us to run queries like:

query.query_path(
    labels=["https://schema.org/performer", 
            "https://schema.org/superEvent", 
            "https://schema.org/organizer"],
    directions=[False, True, True],
    filters=[FilterByType("schema:Person")]
)

During traversal, the method picks outgoing or incoming edges based on the direction:

edges = last.outgoing if forward else last.incoming

For each candidate node, apply_filters checks all filters before including the node in the path:

node_passes_filters = self.apply_filters(next_node, filters)
if node_passes_filters:
    new_paths.append(path + [next_node])



--- NOT FOR REPORT ---

Possible functionalities to implement:
    Path & Traversal Queries
    -Find all neighbors of a node (1 hop out, or N hops out)
    -Find all paths between two nodes

    Lookup & Filter Queries
    -Count how many edges of a given label exist in the graph
    -List all unique edge labels in the graph

    Validation
    -Verify referential integrity — every edge target actually exists as a proper node
    -Check required properties — does every Event node have a schema:name?
    -Detect cycles — is there a loop in the graph?

\section{Task 10}

\section{Task 11}

\section{Task 12}

\section{Task 13}

\section{Task 14}

\section{Conclusion}
Summarize the work completed and key findings.

\end{document}
